\begin{table}[H]
\caption{\raggedright Linimasa Pengerjaan Penelitian}
\label{tab:timeline_overview}
\footnotesize
\begin{tabular}{|>{\raggedright\arraybackslash}m{3.25cm}|>{\centering\arraybackslash}m{2.5cm}|>{\centering\arraybackslash}m{1.25cm}|>{\raggedright\arraybackslash}m{5cm}|}
\hline
\multicolumn{1}{|>{\centering\arraybackslash}m{3.25cm}|}{\textbf{Fase}} & \textbf{Periode} & \textbf{Durasi} & \multicolumn{1}{>{\centering\arraybackslash}m{5cm}|}{\textbf{Luaran Utama}} \\
\hline
Fase 1: Persiapan dan Analisis & Oktober - November 2025 & 2 bulan & Dokumen kebutuhan, desain infrastruktur, dan studi literatur yang terstruktur \\
\hline
Fase 2: Implementasi Infrastruktur & Desember 2025 & 1 bulan & Klaster Kubernetes, lapisan penyimpanan, dan penyiapan \textit{GitOps} \\
\hline
Fase 3: Implementasi Layanan Fitur & Januari 2026 & 1 bulan & Penerapan Feast, fitur sampel, dan \textit{pipeline} materialisasi awal \\
\hline
Fase 4: Implementasi Siklus Hidup Model & Februari 2026 & 1 bulan & Penerapan Kubeflow, penyiapan MLflow, dan alur kerja pelatihan model \\
\hline
Fase 5: Penyajian \& Fitur Lanjutan & Maret 2026 & 1 bulan & Penerapan KServe, deteksi \textit{drift}, serta pemrosesan data lanjutan \\
\hline
Fase 6: Integrasi dan Pengujian & April 2026 & 1 bulan & Penerapan DataHub, integrasi \textit{end-to-end}, dan pengujian dengan pengguna \\
\hline
Fase 7: Dokumentasi dan Finalisasi & Mei 2026 & 1 bulan & Penyusunan tesis, persiapan presentasi, dan penyelesaian revisi akhir \\
\hline
\end{tabular}
\end{table}